% Emacs settings: -*-mode: latex; TeX-master: "manual.tex"; -*-

\chapter{Source components}
\label{c:source}
\index{Sources|textbf}
\index{Library!Components!sources}

\MCS contains a number of different source components,
and any simulation will usually contain exactly one of these sources.
The main function of a source is to determine a set of initial
parameters $(\textbf{r}, \textbf{v}, t)$
for each neutron ray. This is done by Monte Carlo choices from
suitable distributions. For example, in most present sources
the initial position is
found from a uniform distribution over the source surface,
which can be chosen to be either circular or rectangular.
The initial neutron velocity is selected within an interval
of either the corresponding energy or the corresponding wavelength.
Polarization is not relevant for sources,
and we initialize the neutron average spin to zero: $\textbf{s}=(0,0,0)$.

For time-of-flight sources, the choice of the emission time, $t$,
is being made on basis of detailed analytical expressions.
For other sources, $t$ is set to zero.
In the case one would like to use a steady state source
with time-of-flight settings,
the emission time of each neutron ray should be determined using
a Monte Carlo choice. This may be achieved by
the \verb+EXTEND+ keyword in the instrument description source
as in the example below:\index{Keyword!EXTEND}

\begin{lstlisting}
  TRACE

  COMPONENT MySource=Source_gen(...) AT (...)
  EXTEND
  %{
    t = 1e-3*randpm1(); /* set time to +/- 1 ms */
  %}
\end{lstlisting}

\subsection{Neutron flux}
\label{s:neutron-flux}
The flux of the sources deserves special attention. The total neutron
intensity is defined as the sum of weights of all emitted neutron rays
during one simulation
(the unit of total neutron weight is thus neutrons per second).
The flux, $\psi$, at an instrument is defined as intensity per area perpendicular
to the beam direction.

The source flux, $\Phi$, is defined in different units:
the number of neutrons emitted per second from a
1~cm$^2$ area on the source surface,
with direction within a 1~ster.\ solid angle,
and with wavelength within a 1 {\AA} interval.
The total intensity of real neutrons emitted towards a given diaphragm
(units: n/sec) is therefore (for constant $\Phi$):
\begin{equation}
I_\textrm{total} = \Phi A \Delta\Omega \Delta\lambda ,
\end{equation}
where $A$ is the source area, $\Delta\Omega$ is the solid angle of the
diaphragm as seen from the source surface, and $\Delta\lambda$ is the
width of the wavelength interval in which neutrons are emitted (assuming
a uniform wavelength spectrum).

The simulations are performed so that detector intensities
are independent of the number of neutron histories simulated
(although more neutron histories will give better statistics).
If $N_\textrm{sim}$ denotes the number of
neutron histories to simulate, the initial neutron weight $p_0$ must be set to
\begin{equation}
\label{proprule}
p_0 = \frac{N_\textrm{total}}{N_\textrm{sim}} =
    \frac{\Phi(\lambda)}{N_\textrm{sim}} A \Omega \Delta\lambda ,
\end{equation}
where the source flux is now given a $\lambda$-dependence.

As a start, we recommend new \MCS users to use the
\textbf{Source\_simple} component.
Slightly more realistic sources are \textbf{Source\_Maxwell\_3} for
continuous sources or \textbf{Moderator} for time-of-flight sources.

Optimizers can dramatically improve the statistics, but may occasionally
give wrong results, due to misleaded optimization.
You should always check such simulations with (shorter) non-optimized ones.

Other ways to speed-up simulations are to read events from a file.
See section \ref{sources-seealso} for details.

\begin{figure}
  \begin{center}
    \includegraphics[width=0.75\textwidth]{figures/sources}
  \end{center}
\caption{A circular source component (at z=0) emitting neutron events randomly, either from a model, or from a data file.}
\label{f:source}
\end{figure}

\newpage
\section{The \texttt{Source\_simple} McStas Component}
A circular neutron source with flat energy spectrum and arbitrary flux

\subsection*{Identification}
\begin{itemize}
  \item \textbf{Site:} 
  \item \textbf{Author:} Kim Lefmann
  \item \textbf{Origin:} Risoe
  \item \textbf{Date:} October 30, 1997
\end{itemize}

\subsection*{Description}
\begin{lstlisting}
The routine is a circular neutron source, which aims at a square target
centered at the beam (in order to improve MC-acceptance rate).  The angular
divergence is then given by the dimensions of the target.
The neutron energy is uniformly distributed between lambda0-dlambda and
lambda0+dlambda or between E0-dE and E0+dE.
The flux unit is specified in n/cm2/s/st/energy unit (meV or Angs).

This component replaces Source_flat, Source_flat_lambda,
Source_flux and Source_flux_lambda.

Example: Source_simple(radius=0.1, dist=2, focus_xw=.1, focus_yh=.1, E0=14, dE=2)
\end{lstlisting}

\subsection*{Input parameters}
Parameters in \textbf{boldface} are required; the others are optional.

\begin{longtable}{p{0.22\textwidth}p{0.12\textwidth}p{0.46\textwidth}p{0.14\textwidth}}
\toprule
\textbf{Name} & \textbf{Unit} & \textbf{Description} & \textbf{Default} \\
\midrule
\endhead
radius & m & Radius of circle in (x,y,0) plane where neutrons are generated. & 0.1 \\
yheight & m & Height of rectangle in (x,y,0) plane where neutrons are generated. & 0 \\
xwidth & m & Width of rectangle in (x,y,0) plane where neutrons are generated. & 0 \\
dist & m & Distance to target along z axis. & 0 \\
focus\_xw & m & Width of target & .045 \\
focus\_yh & m & Height of target & .12 \\
E0 & meV & Mean energy of neutrons. & 0 \\
dE & meV & Energy half spread of neutrons (flat or gaussian sigma). & 0 \\
lambda0 & AA & Mean wavelength of neutrons. & 0 \\
dlambda & AA & Wavelength half spread of neutrons. & 0 \\
flux & 1/(s*cm**2*st*energy unit) & flux per energy unit, Angs or meV if flux=0, the source emits 1 in 4*PI whole space. & 1 \\
gauss & 1 & Gaussian (1) or Flat (0) energy/wavelength distribution & 0 \\
target\_index & 1 & relative index of component to focus at, e.g. next is +1 this is used to compute 'dist' automatically. & 1 \\
\bottomrule
\end{longtable}

\subsection*{Links}
\begin{itemize}
  \item \href{run:/home/willend/willend-McCode/mcstas-comps/sources/Source_simple.comp}{Source code} for \texttt{Source\_simple.comp}.
\end{itemize}
\IfFileExists{Source_simple_static.tex}{\input{Source_simple_static.tex}}{}

\section{The \texttt{Source\_div} McStas Component}
Neutron source with Gaussian or uniform divergence

\subsection*{Identification}
\begin{itemize}
  \item \textbf{Site:} 
  \item \textbf{Author:} KL
  \item \textbf{Origin:} Risoe
  \item \textbf{Date:} November 20, 1998
\end{itemize}

\subsection*{Description}
\begin{lstlisting}
The routine is a rectangular neutron source, which has a gaussian or uniform
divergent output in the forward direction.
The neutron energy is distributed between lambda0-dlambda and
lambda0+dlambda or between E0-dE and E0+dE. The flux unit is specified
in n/cm2/s/st/energy unit (meV or Angs).
In the case of uniform distribution (gauss=0), angles are uniformly distributed
between -focus_aw and +focus_aw as well as -focus_ah and +focus_ah.
For Gaussian distribution (gauss=1), 'focus_aw' and 'focus_ah' define the
FWHM of a Gaussian distribution. Energy/wavelength distribution is also
Gaussian.

Example: Source_div(xwidth=0.1, yheight=0.1, focus_aw=2, focus_ah=2, E0=14, dE=2, gauss=0)

%VALIDATION
Feb 2005: tested by Kim Lefmann    (o.k.)
Apr 2005: energy distribution used in external tests of Fermi choppers (o.k.)
Jun 2005: wavelength distribution used in external tests of velocity selectors (o.k.)
Validated by: K. Lieutenant

%BUGS
distribution is uniform in (hor. and vert.) angle (relative to moderator normal),
therefore not suited for large angles
\end{lstlisting}

\subsection*{Input parameters}
Parameters in \textbf{boldface} are required; the others are optional.

\begin{longtable}{p{0.22\textwidth}p{0.12\textwidth}p{0.46\textwidth}p{0.14\textwidth}}
\toprule
\textbf{Name} & \textbf{Unit} & \textbf{Description} & \textbf{Default} \\
\midrule
\endhead
\textbf{xwidth} & m & Width of source &  \\
\textbf{yheight} & m & Height of source &  \\
\textbf{focus\_aw} & deg & FWHM (Gaussian) or maximal (uniform) horz. width divergence &  \\
\textbf{focus\_ah} & deg & FWHM (Gaussian) or maximal (uniform) vert. height divergence &  \\
E0 & meV & Mean energy of neutrons. & 0.0 \\
dE & meV & Energy half spread of neutrons. & 0.0 \\
lambda0 & Ang & Mean wavelength of neutrons (only relevant for E0=0) & 0.0 \\
dlambda & Ang & Wavelength half spread of neutrons. & 0.0 \\
gauss & 0|1 & Criterion: 0: uniform, 1: Gaussian distributions & 0 \\
flux & 1/(s cm 2 st energy\_unit) & flux per energy unit, Angs or meV & 1 \\
\bottomrule
\end{longtable}

\subsection*{Links}
\begin{itemize}
  \item \href{run:/home/willend/willend-McCode/mcstas-comps/sources/Source_div.comp}{Source code} for \texttt{Source\_div.comp}.
\end{itemize}
\IfFileExists{Source_div_static.tex}{\input{Source_div_static.tex}}{}

\section{The \texttt{Source\_Maxwell\_3} McStas Component}
Source with up to three Maxwellian distributions

\subsection*{Identification}
\begin{itemize}
  \item \textbf{Site:} 
  \item \textbf{Author:} Kim Lefmann
  \item \textbf{Origin:} Risoe
  \item \textbf{Date:} March 2001
\end{itemize}

\subsection*{Description}
\begin{lstlisting}
A parametrised continuous source for modelling a (cubic) source
with (up to) 3 Maxwellian distributions.
The source produces a continuous spectrum.
The sampling of the neutrons is uniform in wavelength.

Units of flux: neutrons/cm^2/second/ster
(McStas units are in general neutrons/second)

Example:  PSI cold source T1=150.42 K / 2.51 AA     I1 = 3.67 E11
T2=38.74 K / 4.95 AA      I2 = 3.64 E11
T3=14.84 K / 9.5 AA       I3 = 0.95 E11
\end{lstlisting}

\subsection*{Input parameters}
Parameters in \textbf{boldface} are required; the others are optional.

\begin{longtable}{p{0.22\textwidth}p{0.12\textwidth}p{0.46\textwidth}p{0.14\textwidth}}
\toprule
\textbf{Name} & \textbf{Unit} & \textbf{Description} & \textbf{Default} \\
\midrule
\endhead
size & m & Edge of cube shaped source (for backward compatibility) & 0 \\
yheight & m & Height of rectangular source & 0 \\
xwidth & m & Width of rectangular source & 0 \\
\textbf{Lmin} & AA & Lower edge of lambda distribution &  \\
\textbf{Lmax} & AA & Upper edge of lambda distribution &  \\
\textbf{dist} & m & Distance from source to focusing rectangle; at (0,0,dist) &  \\
\textbf{focus\_xw} & m & Width of focusing rectangle &  \\
\textbf{focus\_yh} & m & Height of focusing rectangle &  \\
\textbf{T1} & K & 1st temperature of thermal distribution &  \\
T2 & K & 2nd temperature of thermal distribution & 300 \\
T3 & K & 3nd temperature of  - - - & 300 \\
\textbf{I1} & 1/(cm**2*st) & flux, 1 (in flux units, see above) &  \\
I2 & 1/(cm**2*st) & flux, 2 (in flux units, see above) & 0 \\
I3 & 1/(cm**2*st) & flux, 3  - - - & 0 \\
target\_index & 1 & relative index of component to focus at, e.g. next is +1 this is used to compute 'dist' automatically. & 1 \\
lambda0 & AA & Mean wavelength of neutrons. & 0 \\
dlambda & AA & Wavelength spread of neutrons. & 0 \\
\bottomrule
\end{longtable}

\subsection*{Links}
\begin{itemize}
  \item \href{run:/home/willend/willend-McCode/mcstas-comps/sources/Source_Maxwell_3.comp}{Source code} for \texttt{Source\_Maxwell\_3.comp}.
\end{itemize}
\IfFileExists{Source_Maxwell_3_static.tex}{\input{Source_Maxwell_3_static.tex}}{}

\input{sources/Source_gen}

\newpage
\section{The \texttt{Moderator} McStas Component}
A simple pulsed source for time-of-flight.

\subsection*{Identification}
\begin{itemize}
  \item \textbf{Site:} 
  \item \textbf{Author:} KN, M.Hagen
  \item \textbf{Origin:} Risoe
  \item \textbf{Date:} August 1998
\end{itemize}

\subsection*{Description}
\begin{lstlisting}
Produces a simple time-of-flight spectrum, with a flat energy distribution

Example: Moderator(radius = 0.0707, dist = 9.035, focus_xw = 0.021, focus_yh = 0.021, Emin = 10, Emax = 15, Ec = 9.0, t0 = 37.15, gamma = 39.1)
\end{lstlisting}

\subsection*{Input parameters}
Parameters in \textbf{boldface} are required; the others are optional.

\begin{longtable}{p{0.22\textwidth}p{0.12\textwidth}p{0.46\textwidth}p{0.14\textwidth}}
\toprule
\textbf{Name} & \textbf{Unit} & \textbf{Description} & \textbf{Default} \\
\midrule
\endhead
radius & m & Radius of source & 0.07 \\
\textbf{Emin} & meV & Lower edge of energy distribution &  \\
\textbf{Emax} & meV & Upper edge of energy distribution &  \\
dist & m & Distance from source to the focusing rectangle & 0 \\
focus\_xw & m & Width of focusing rectangle & 0.02 \\
focus\_yh & m & Height of focusing rectangle & 0.02 \\
t0 & mus & decay constant for low-energy neutrons & 37.15 \\
Ec & meV & Critical energy, below which the flux decay is constant & 9.0 \\
gamma & meV & energy dependence of decay time & 39.1 \\
target\_index & 1 & relative index of component to focus at, e.g. next is +1 this is used to compute 'dist' automatically. & 1 \\
flux & 1/(s cm 2 st meV) & flux & 1 \\
\bottomrule
\end{longtable}

\subsection*{Links}
\begin{itemize}
  \item \href{run:/home/willend/willend-McCode/mcstas-comps/sources/Moderator.comp}{Source code} for \texttt{Moderator.comp}.
\end{itemize}
\IfFileExists{Moderator_static.tex}{\input{Moderator_static.tex}}{}

%\section{The \texttt{ISIS\_moderator} McStas Component}
ISIS Moderators

\subsection*{Identification}
\begin{itemize}
  \item \textbf{Site:} 
  \item \textbf{Author:} S. Ansell and D. Champion
  \item \textbf{Origin:} ISIS
  \item \textbf{Date:} AUGUST 2005
\end{itemize}

\subsection*{Description}
\begin{lstlisting}
Produces a TS1 or TS2 ISIS moderator distribution.  The Face argument determines which moderator
is to be sampled. Each face uses a different Etable file (the location of which is determined by
the MCTABLES environment variable).  Neutrons are created having a range of energies
determined by the Emin and Emax arguments.  Trajectories are produced such that they pass
through the moderator face (defined by modXsixe and yheight) and a focusing rectangle
(defined by xh,focus_yh and dist).  Please download the documentation for precise instructions for use.

Example:   ISIS_moderator(Face ="water", Emin = 49.0,Emax = 51.0, dist = 1.0, focus_xw = 0.01,
focus_yh = 0.01, xwidth = 0.074, yheight = 0.074, CAngle = 0.0,SAC= 1)
\end{lstlisting}

\subsection*{Input parameters}
Parameters in \textbf{boldface} are required; the others are optional.

\begin{longtable}{p{0.22\textwidth}p{0.12\textwidth}p{0.46\textwidth}p{0.14\textwidth}}
\toprule
\textbf{Name} & \textbf{Unit} & \textbf{Description} & \textbf{Default} \\
\midrule
\endhead
Face & word & Name of the face - TS2=groove,hydrogen,narrow,broad - & "hydrogen" \\
Emin & meV & Lower edge of energy distribution & 49.0 \\
Emax & meV & Upper edge of energy distribution & 51.0 \\
dist & m & Distance from source to the focusing rectangle & 1.0 \\
focus\_xw & m & Width of focusing rectangle & 0.01 \\
focus\_yh & m & Height of focusing rectangle & 0.01 \\
xwidth & m & Moderator vertical size & 0.074 \\
yheight & m & Moderator Horizontal size & 0.074 \\
CAngle & radians & Angle from the centre line & 0.0 \\
SAC & boolean & Apply solid angle correction or not (1/0) & 1 \\
Lmin & meV & Lower edge of wavelength distribution & 0 \\
Lmax & meV & Upper edge of wavelength distribution & 0 \\
target\_index & 1 & relative index of component to focus at, e.g. next is +1 & 1 \\
verbose & boolean & Echo status information at everu 1e5'th neutron & 0 \\
\bottomrule
\end{longtable}

\subsection*{Links}
\begin{itemize}
  \item \href{run:/home/willend/willend-McCode/mcstas-comps/contrib/ISIS_moderator.comp}{Source code} for \texttt{ISIS\_moderator.comp}.
  \item Further \textless{}A HREF="http://www.isis.rl.ac.uk/Computing/Software/MC/index.htm"\textgreater{}information\textless{}/A\textgreater{} should be
  \item downloaded from the ISIS MC website.
  \item *******************************************************************************/
\end{itemize}
\IfFileExists{ISIS_moderator_static.tex}{\input{ISIS_moderator_static.tex}}{}

\newpage
\section{The \texttt{Source\_adapt} McStas Component}
Neutron source with adaptive importance sampling

\subsection*{Identification}
\begin{itemize}
  \item \textbf{Site:} 
  \item \textbf{Author:} Kristian Nielsen
  \item \textbf{Origin:} Risoe
  \item \textbf{Date:} 1999
\end{itemize}

\subsection*{Description}
\begin{lstlisting}
Rectangular source with flat energy or wavelength distribution that
uses adaptive importance sampling to improve simulation efficiency.
Works together with the Adapt_check component.

The source divides the three-dimensional phase space of (energy,
horizontal position, horizontal divergence) into a number of
rectangular bins. The probability for selecting neutrons from each
bin is adjusted so that neutrons that reach the Adapt_check
component with high weights are emitted more frequently than those
with low weights. The adjustment is made so as to attemt to make
the weights at the Adapt_check components equal.

Focusing is achieved by only emitting neutrons towards a rectangle
perpendicular to and placed at a certain distance along the Z axis.
Focusing is only approximate (for simplicity); neutrons are also
emitted to pass slightly above and below the focusing rectangle,
more so for wider focusing.

In order to prevent false learning, a parameter beta sets a
fraction of the neutrons that are emitted uniformly, without regard
to the adaptive distribution. The parameter alpha sets an initial
fraction of neutrons that are emitted with low weights; this is
done to prevent early neutrons with rare initial parameters but
high weight to ruin the statistics before the component adapts its
distribution to the problem at hand. Good general-purpose values
for these parameters are alpha = beta = 0.25.

%VALIDATION
This component is not validated. It does not work properly with MPI.
\end{lstlisting}

\subsection*{Input parameters}
Parameters in \textbf{boldface} are required; the others are optional.

\begin{longtable}{p{0.22\textwidth}p{0.12\textwidth}p{0.46\textwidth}p{0.14\textwidth}}
\toprule
\textbf{Name} & \textbf{Unit} & \textbf{Description} & \textbf{Default} \\
\midrule
\endhead
N\_E & 1 & Number of bins in energy (or wavelength) dimension & 20 \\
N\_xpos & 1 & Number of bins in horizontal position & 20 \\
N\_xdiv & 1 & Number of bins in horizontal divergence & 20 \\
xmin & m & Left edge of rectangular source & 0 \\
xmax & m & Right edge & 0 \\
ymin & m & Lower edge & 0 \\
ymax & m & Upper edge & 0 \\
xwidth & m & Width of source & 0 \\
yheight & m & Height of source & 0 \\
filename & string & Optional filename for adaptive distribution output & 0 \\
dist & m & Distance to target rectangle along z axis & 0 \\
focus\_xw & m & Width of target & 0.05 \\
focus\_yh & m & Height of target & 0.1 \\
E0 & meV & Mean energy of neutrons & 0 \\
dE & meV & Energy spread (energy range is from E0-dE to E0+dE) & 0 \\
lambda0 & AA & Mean wavelength of neutrons (if energy not specified) & 0 \\
dlambda & AA & Wavelength spread half width & 0 \\
flux &  & (1/(cm 2 AA st)) Absolute source flux & 1e13 \\
target\_index & 1 & relative index of component to focus at, e.g. next is +1 this is used to compute 'dist' automatically. & 1 \\
alpha & 1 & Learning cut-off factor (0 \textless{} alpha \textless{}= 1) & 0.25 \\
beta & 1 & Aggressiveness of adaptive algorithm (0 \textless{} beta \textless{}= 1) & 0.25 \\
\bottomrule
\end{longtable}

\subsection*{Links}
\begin{itemize}
  \item \href{run:/home/willend/willend-McCode/mcstas-comps/sources/Source_adapt.comp}{Source code} for \texttt{Source\_adapt.comp}.
\end{itemize}
\IfFileExists{Source_adapt_static.tex}{\input{Source_adapt_static.tex}}{}

\section{The \texttt{Adapt\_check} McStas Component}
Optimization specifier for the Source\_adapt component.

\subsection*{Identification}
\begin{itemize}
  \item \textbf{Site:} 
  \item \textbf{Author:} Kristian Nielsen
  \item \textbf{Origin:} Risoe
  \item \textbf{Date:} 1999
\end{itemize}

\subsection*{Description}
\begin{lstlisting}
This components works together with the Source_adapt component, and
is used to define the criteria for selecting which neutrons are
considered "good" in the adaptive algorithm. The name of the
associated Source_adapt component in the instrument definition is
given as parameter. The component is special in that its position
does not matter; all neutrons that have not been absorbed prior to
the component are considered "good".

Example: Adapt_check(source_comp="MySource")
\end{lstlisting}

\subsection*{Input parameters}
Parameters in \textbf{boldface} are required; the others are optional.

\begin{longtable}{p{0.22\textwidth}p{0.12\textwidth}p{0.46\textwidth}p{0.14\textwidth}}
\toprule
\textbf{Name} & \textbf{Unit} & \textbf{Description} & \textbf{Default} \\
\midrule
\endhead
\textbf{source\_comp} & string & The name of the Source\_adapt component in the instrument definition. &  \\
\bottomrule
\end{longtable}

\subsection*{Links}
\begin{itemize}
  \item \href{run:/home/willend/willend-McCode/mcstas-comps/sources/Adapt_check.comp}{Source code} for \texttt{Adapt\_check.comp}.
\end{itemize}
\IfFileExists{Adapt_check_static.tex}{\input{Adapt_check_static.tex}}{}

\newpage
\section{The \texttt{Source\_Optimizer} McStas Component}
A component that optimizes the neutron flux passing through the
Source\_Optimizer in order to have the maximum flux at the
\textless{}b\textgreater{}Monitor\_Optimizer\textless{}/b\textgreater{} position.

\subsection*{Identification}
\begin{itemize}
  \item \textbf{Site:} 
  \item \textbf{Author:} \textless{}a href="mailto:farhi@ill.fr"\textgreater{}Emmanuel Farhi\textless{}/a\textgreater{}
  \item \textbf{Origin:} \textless{}a href="http://www.ill.fr"\textgreater{}ILL (France)\textless{}/a\textgreater{}
  \item \textbf{Date:} 17 Sept 1999
\end{itemize}

\subsection*{Description}
\begin{lstlisting}
Principle: The optimizer first (step 1) computes neutron state parameter
limits passing in the Source_Optimizer, and then (step 2) records a Reference
source as well as the state (at Source_Optimizer position) of neutrons
reaching Monitor. The optimized source is defined as a fraction of the
Reference source plus the distribution of 'good' neutrons reaching the
Monitor. The optimization then starts (step 3), and focuses new neutrons on
the Monitor_Optimizer. In fact it changes 'bad' neutrons into 'good' ones
(that reach the Monitor), acting on their position, spin and divergence or
velocity. The overall Monitor flux is kept during process. The energy and
polarisation distributions are kept during optimization as far as possible
during optimisation. The optimization method considers that all neutron
parameters - (x,y), (vx,vy,vz) or (vx/v2,vy/v2,v2), (sx,sy,sz) or
(sx/s2,sy/s2,s2) - are independent.

Options: The optimized source can be computed regularly ('continuous'
option) or only once ('not continuous'). The time spent in steps 1 and 2 can
be reduced for a shorter optimization ('auto').  The neutrons passing during
steps 1 and 2 can be smoothed for a better neutron weight distribution
('smooth' option).

Source_optimizer can be placed at any position where you want to act on the
flux, for instance just after the source.
Monitor_Optimizer should be placed at position(s) to optimize.
I prefer to put one just before the sample.

Default parameters bins, step, and keep are 10, 10%  and 10% respectively.
The option string can be empty (""), which stands for default configuration
that works fine in usual cases:

options="continuous optimization, auto mode, smooth, SetXY+SetDivV+SetDivS"

<b>Possible options are</b>
continuous      for continuous source optimization (default).
verbose         displays optimization process (debug purpose).
auto            uses the shortest possible 'step 1' and 'step 2' and sets 'step' value as required (default).
smooth          remove possible spikes generated in steps 1 and 2 (default is smooth).
inactivate      to inactivate the Optimizer.
no or not       revert next option
bins=[value=10] set the Number of cells for sampling neutron states
step=[value=10] Optimizer step in % of simulation.
keep=[value=10] Percentage of initial source distribution that is kept
file=[name]     Filename where to save optimized source distributions (no file is generated if not given. Default ext. is .src)
SetXY           Keywords to indicate what may be changed during
SetV            optimisation. Default is position, divergence and spin
SetS            direction ("SetXY+SetDivV+SetdivS"). Choosing the speed
SetDivV         or spin optimization (SetV or SetS) may modify the energy
SetDivS         or polarisation distribution (norm of V and S) as the three components are then independent.

Parameters bins, step and keep can also be entered as optional parameters.

<b>EXAMPLE</b>: I use the following settings

optim_s = Source_Optimizer(options="please be clever") (same as empty)
(...)
Monitor_Optimizer(xmin=-0.05, xmax=0.05, ymin=-0.05, ymax=0.05,
optim_comp = "optim_s")

A good optimization needs to record enough non optimized neutrons on Monitor
during step 2. Typical enhancement in computation speed is by a factor 20.
This component usually works well.

<b>NOTE:</b> You must be aware that in some cases (SetV and SetS),
the optimization might sligtly afect the energy or spin distribution of the
source. The optimizer tries to do its best anyway.
Also, some 'spikes' may sometime appear in monitor signals in the course of
the optimization, coming from non-optimized neutrons with original weight.
The 'smooth' option minimises this effect (on by default).
\end{lstlisting}

\subsection*{Input parameters}
Parameters in \textbf{boldface} are required; the others are optional.

\begin{longtable}{p{0.22\textwidth}p{0.12\textwidth}p{0.46\textwidth}p{0.14\textwidth}}
\toprule
\textbf{Name} & \textbf{Unit} & \textbf{Description} & \textbf{Default} \\
\midrule
\endhead
bins & 1 & Number of cells for sampling neutron states. & 10 \\
step & 0-100 & Optimizer step in percent of simulation. & 0.1 \\
keep & 0-100 & Percentage of initial source distribution that is kept. & 0.1 \\
options & str & string of options. See \textless{}b\textgreater{}Description\textless{}b\textgreater{} & 0 \\
\bottomrule
\end{longtable}

\subsection*{Links}
\begin{itemize}
  \item \href{run:/home/willend/willend-McCode/mcstas-comps/sources/Source_Optimizer.comp}{Source code} for \texttt{Source\_Optimizer.comp}.
\end{itemize}
\IfFileExists{Source_Optimizer_static.tex}{\input{Source_Optimizer_static.tex}}{}

\section{The \texttt{Monitor\_Optimizer} McStas Component}
To be used after the \textless{}b\textgreater{}Source\_Optimizer\textless{}/b\textgreater{} component

\subsection*{Identification}
\begin{itemize}
  \item \textbf{Site:} 
  \item \textbf{Author:} \textless{}a href="mailto:farhi@ill.fr"\textgreater{}Emmanuel Farhi\textless{}/a\textgreater{}
  \item \textbf{Origin:} \textless{}a href="http://www.ill.fr"\textgreater{}ILL (France)\textless{}/a\textgreater{}
  \item \textbf{Date:} 17 Sept 1999
\end{itemize}

\subsection*{Description}
\begin{lstlisting}
A component that optimizes the neutron flux passing through the
<b>Source_Optimizer</b> in order to have the maximum flux at the
Monitor_Optimizer position(s).
<b>Source_optimizer</b> should be placed just after the source.
Monitor_Optimizer should be placed at the position to optimize.
I prefer to put one just before the sample.

See <a href="Source_Optimizer.html">Source_Optimizer</a> for
usage example and additional informations.
\end{lstlisting}

\subsection*{Input parameters}
Parameters in \textbf{boldface} are required; the others are optional.

\begin{longtable}{p{0.22\textwidth}p{0.12\textwidth}p{0.46\textwidth}p{0.14\textwidth}}
\toprule
\textbf{Name} & \textbf{Unit} & \textbf{Description} & \textbf{Default} \\
\midrule
\endhead
\textbf{optim\_comp} & str & name of the Source\_Optimizer component in the instrument definition. Do not use quotes (no quotes) &  \\
xmin & m & Lower x bound of monitor opening & -0.1 \\
xmax & m & Upper x bound of monitor opening & 0.1 \\
ymin & m & Lower y bound of monitor opening & -0.1 \\
ymax & m & Upper y bound of monitor opening & 0.1 \\
xwidth & m & Width of monitor. Overrides xmin,xmax. & 0 \\
yheight & m & Height of monitor. Overrides ymin,ymax. & 0 \\
\bottomrule
\end{longtable}

\subsection*{Links}
\begin{itemize}
  \item \href{run:/home/willend/willend-McCode/mcstas-comps/sources/Monitor_Optimizer.comp}{Source code} for \texttt{Monitor\_Optimizer.comp}.
  \item \textless{}a href="Source\_Optimizer.html"\textgreater{}Source\_Optimizer\textless{}/a\textgreater{}
\end{itemize}
\IfFileExists{Monitor_Optimizer_static.tex}{\input{Monitor_Optimizer_static.tex}}{}

\newpage
\section{Other sources components: contributed pulsed sources, virtual sources (event files)}
\label{sources-seealso}

There are many other source definitions in \MCS .

Detailed pulsed source components are available for new facilities
in a number of contributed components:
\begin{itemize}
\item SNS (\textbf{contrib/SNS\_source}),
\item ISIS (\textbf{contrib/ISIS\_moderator}),
\item ESS-project (\textbf{sources/ESS\_butterfly}), the current, detailed
  time-of-flight moderator/target-station model for the European Spallation
  Source, superseding the earlier \textbf{ESS\_moderator\_short} (now in
  \textbf{obsolete}).
\end{itemize}

When no analytical model (e.g. a Maxwellian distribution) exists,
one may have access to measurements, estimated flux distributions, or
neutron event files. As of \MCS 3.x, the recommended, portable way to read
and write such neutron event lists is the MCPL format (see the
\textbf{misc/MCPL\_input} and \textbf{misc/MCPL\_output} components,
section~\ref{s:mcpl}), which has backends for \MCS/McXtrace as well as
several other Monte Carlo transport codes (MCNP, PHITS, Geant4, Vitess,
SIMRES). The older, code-specific event-file components
(\textbf{Virtual\_input}/\textbf{Virtual\_output},
\textbf{Virtual\_mcnp\_input}/\textbf{Virtual\_mcnp\_output},
\textbf{Virtual\_tripoli4\_input}/\textbf{Virtual\_tripoli4\_output},
\textbf{Vitess\_input}/\textbf{Vitess\_output}) still exist for backward
compatibility but have moved to the \textbf{obsolete} component category
(see the Component Manual's Obsolete chapter); new instrument work should
prefer MCPL.\index{Vitess} \index{Tripoli} \index{MCNP}

Finally, \textbf{optics/Filter\_gen} reads a 1D distribution from a file,
and may either modify or set the flux according to it. See section
\ref{filter-gen}.
